% GENERATED FILE. Edit canonical lessons, then run npm run generate:books.
% canonical-source-hash: fnv1a64:dddfc4cd00d2d1c3
% canonical-lessons: JA-C511-hikiokosu, JA-C511-seika, JA-C511-shoujiru, JA-C511-kikime, JA-C511-moto

\chapter{What Caused It}
\label{ch:ja-what-caused-it}
\hlchaptermodality{511}

\begin{chapteropening}
\textbf{By the end of this chapter:} \emph{I can talk about causes and effects.}

Complete the last lesson of chapter 511: I can talk about causes and effects.
\end{chapteropening}

\section[\texorpdfstring{\ja{ひきおこす} (hikiokosu)}{ひきおこす (hikiokosu)}]{\ja{ひきおこす} (hikiokosu) — to cause; to bring about}

\label{lesson:JA-C511-hikiokosu}



\begin{quote}
Before the new one: say the Japanese for to make good use of, then the Japanese for to be blessed (with).
\end{quote}

\subsection*{You'll want to know: \ja{ひきおこす}}
\textbf{\ja{ひきおこす}} — \emph{hikiokosu} — \textquotedblleft{}to cause; to bring about\textquotedblright{}.

\begin{quote}
\textbf{\ja{もし} \ja{かれが} \ja{あぶない} \ja{うんてんを} \ja{したら、じこを} \ja{ひきおこすでしょう。}} — \emph{moshi kare ga abunai unten o shitara, jiko o hikiokosu deshō} — \textquotedblleft{}If he drives dangerously, he will probably cause an accident.\textquotedblright{}
\end{quote}

Polite \textbf{\ja{ひきおこします}}, \emph{hikiokoshimasu}; a godan verb. Past (polite): \textbf{\ja{ゆきが} \ja{じこを} \ja{ひきおこしました。}}, \emph{yuki ga jiko o hikiokoshimashita}.

\begin{cousinweb}
\textbf{\ja{ひきおこす}} (\emph{hikiokosu}) is a native Japanese word (wago), inherited rather than borrowed from Chinese, built from \emph{hiku}, \textquotedblleft{}to pull\textquotedblright{}, and \emph{okosu}, \textquotedblleft{}to raise\textquotedblright{}; its kanji was attached later as a kun'yomi (JMdict, KANJIDIC2).
\end{cousinweb}

\begin{grammarlens}[title={What you've built}]
One more word for saying what happens if.
\end{grammarlens}

\subsection*{Your turn}
\begin{itemize}
\raggedright
  \item \emph{Say it:} \emph{hikiokosu}
  \item \emph{Say it:} \emph{hikiokosu}, once more
  \item \emph{Say it:} \emph{moshi kare ga abunai unten o shitara, jiko o hikiokosu deshō} — If he drives dangerously, he will probably cause an accident.
  \item \emph{Recall:} say the Japanese for to make good use of, then the Japanese for to be blessed (with), then say \emph{hikiokosu} again
\end{itemize}

\begin{tcolorbox}[breakable,colback=teal!4,colframe=teal!35!black,title={Before you move on}]
What does \ja{ひきおこす} mean? (\textquotedblleft{}To cause; to bring about\textquotedblright{}.) Say it once more.
\end{tcolorbox}
\section[\texorpdfstring{\ja{せいか} (seika)}{せいか (seika)}]{\ja{せいか} (seika) — an achievement; the fruits (of effort)}

\label{lesson:JA-C511-seika}



\begin{quote}
Before the new one: say the Japanese for to be blessed (with), then the Japanese for to cause.
\end{quote}

\subsection*{You'll want to know: \ja{せいか}}
\textbf{\ja{せいか}} — \emph{seika} — \textquotedblleft{}an achievement; the fruits (of effort)\textquotedblright{}.

\begin{quote}
\textbf{\ja{もし} \ja{れんしゅうの} \ja{せいかが} \ja{あったら、せんしゅは} \ja{よろこぶでしょう。}} — \emph{moshi renshū no seika ga attara, senshu wa yorokobu deshō} — \textquotedblleft{}If the practice brings results, the players will be glad.\textquotedblright{}
\end{quote}

\begin{cousinweb}
\textbf{\ja{せいか}} (\emph{seika}) is Sino-Japanese (kango): two Chinese characters read with their on'yomi, readings borrowed from Chinese. They are \emph{sei}, \textquotedblleft{}turn into\textquotedblright{} (Middle Chinese \emph{dzyeng}) and \emph{ka}, \textquotedblleft{}fruit\textquotedblright{} (KANJIDIC2; Middle Chinese in Baxter's transcription).
\end{cousinweb}

\begin{grammarlens}[title={What you've built}]
One more word for saying what happens if.
\end{grammarlens}

\subsection*{Your turn}
\begin{itemize}
\raggedright
  \item \emph{Say it:} \emph{seika}
  \item \emph{Say it:} \emph{seika}, once more
  \item \emph{Say it:} \emph{moshi renshū no seika ga attara, senshu wa yorokobu deshō} — If the practice brings results, the players will be glad.
  \item \emph{Recall:} say the Japanese for to be blessed (with), then the Japanese for to cause, then say \emph{seika} again
\end{itemize}

\begin{tcolorbox}[breakable,colback=teal!4,colframe=teal!35!black,title={Before you move on}]
What does \ja{せいか} mean? (\textquotedblleft{}An achievement; the fruits (of effort)\textquotedblright{}.) Say it once more.
\end{tcolorbox}
\section[\texorpdfstring{\ja{しょうじる} (shōjiru)}{しょうじる (shōjiru)}]{\ja{しょうじる} (shōjiru) — to arise; to come about}

\label{lesson:JA-C511-shoujiru}



\begin{quote}
Before the new one: say the Japanese for to cause, then the Japanese for an achievement.
\end{quote}

\subsection*{You'll want to know: \ja{しょうじる}}
\textbf{\ja{しょうじる}} — \emph{shōjiru} — \textquotedblleft{}to arise; to come about\textquotedblright{}.

\begin{quote}
\textbf{\ja{もし} \ja{もんだいが} \ja{しょうじたら、でんわして} \ja{ください。}} — \emph{moshi mondai ga shōjitara, denwa shite kudasai} — \textquotedblleft{}If a problem arises, please call.\textquotedblright{}
\end{quote}

Polite \textbf{\ja{しょうじます}}, \emph{shōjimasu}; an ichidan verb. Past (polite): \textbf{\ja{もんだいが} \ja{しょうじました。}}, \emph{mondai ga shōjimashita}.

\begin{cousinweb}
\textbf{\ja{しょうじる}} (\emph{shōjiru}) mixes the two layers, native and Sino-Japanese: Sino-Japanese \emph{shō}, \textquotedblleft{}life\textquotedblright{} (Middle Chinese \emph{srjaeng}), borrowed from Chinese and the native ending \emph{jiru} (JMdict, KANJIDIC2; Middle Chinese in Baxter's transcription).
\end{cousinweb}

\begin{grammarlens}[title={What you've built}]
One more word for saying what happens if.
\end{grammarlens}

\subsection*{Your turn}
\begin{itemize}
\raggedright
  \item \emph{Say it:} \emph{shōjiru}
  \item \emph{Say it:} \emph{shōjiru}, once more
  \item \emph{Say it:} \emph{moshi mondai ga shōjitara, denwa shite kudasai} — If a problem arises, please call.
  \item \emph{Recall:} say the Japanese for to cause, then the Japanese for an achievement, then say \emph{shōjiru} again
\end{itemize}

\begin{tcolorbox}[breakable,colback=teal!4,colframe=teal!35!black,title={Before you move on}]
What does \ja{しょうじる} mean? (\textquotedblleft{}To arise; to come about\textquotedblright{}.) Say it once more.
\end{tcolorbox}
\section[\texorpdfstring{\ja{ききめ} (kikime)}{ききめ (kikime)}]{\ja{ききめ} (kikime) — how well something works}

\label{lesson:JA-C511-kikime}



\begin{quote}
Before the new one: say the Japanese for an achievement, then the Japanese for to arise.
\end{quote}

\subsection*{You'll want to know: \ja{ききめ}}
\textbf{\ja{ききめ}} — \emph{kikime} — \textquotedblleft{}how well something works\textquotedblright{}.

\begin{quote}
\textbf{\ja{この} \ja{くすりの} \ja{ききめは} \ja{びょうきに} \ja{よります。}} — \emph{kono kusuri no kikime wa byōki ni yorimasu} — \textquotedblleft{}How well this medicine works depends on the illness.\textquotedblright{}
\end{quote}

\begin{cousinweb}
\textbf{\ja{ききめ}} (\emph{kikime}) is a native Japanese compound (wago), inherited rather than borrowed from Chinese: built from \emph{kiki}, \textquotedblleft{}working\textquotedblright{} and \emph{me}, \textquotedblleft{}eye\textquotedblright{} (JMdict, KANJIDIC2).
\end{cousinweb}

\begin{grammarlens}[title={What you've built}]
One more word for saying what happens if.
\end{grammarlens}

\subsection*{Your turn}
\begin{itemize}
\raggedright
  \item \emph{Say it:} \emph{kikime}
  \item \emph{Say it:} \emph{kikime}, once more
  \item \emph{Say it:} \emph{kono kusuri no kikime wa byōki ni yorimasu} — How well this medicine works depends on the illness.
  \item \emph{Recall:} say the Japanese for an achievement, then the Japanese for to arise, then say \emph{kikime} again
\end{itemize}

\begin{tcolorbox}[breakable,colback=teal!4,colframe=teal!35!black,title={Before you move on}]
What does \ja{ききめ} mean? (\textquotedblleft{}How well something works\textquotedblright{}.) Say it once more.
\end{tcolorbox}
\section[\texorpdfstring{\ja{もと} (moto)}{もと (moto)}]{\ja{もと} (moto) — the origin; the cause}

\label{lesson:JA-C511-moto}



\begin{quote}
Before the new one: say the Japanese for to arise, then the Japanese for how well something works.
\end{quote}

\subsection*{You'll want to know: \ja{もと}}
\textbf{\ja{もと}} — \emph{moto} — \textquotedblleft{}the origin; the cause\textquotedblright{}.

\begin{quote}
\textbf{\ja{もし} \ja{けんかの} \ja{もとが} \ja{わかったら、おしえて} \ja{ください。}} — \emph{moshi kenka no moto ga wakattara, oshiete kudasai} — \textquotedblleft{}If you find out what started the quarrel, please tell me.\textquotedblright{}
\end{quote}

\begin{cousinweb}
\textbf{\ja{もと}} (\emph{moto}) is a native Japanese word (wago), inherited rather than borrowed from Chinese; its kanji was attached later as a kun'yomi, and its older origin is not settled (JMdict, KANJIDIC2).
\end{cousinweb}

\begin{grammarlens}[title={What you've built}]
One more word for saying what happens if.
\end{grammarlens}

\subsection*{Your turn}
\begin{itemize}
\raggedright
  \item \emph{Say it:} \emph{moto}
  \item \emph{Say it:} \emph{moto}, once more
  \item \emph{Say it:} \emph{moshi kenka no moto ga wakattara, oshiete kudasai} — If you find out what started the quarrel, please tell me.
  \item \emph{Recall:} say the Japanese for to arise, then the Japanese for how well something works, then say \emph{moto} again
\end{itemize}

\begin{tcolorbox}[breakable,colback=teal!4,colframe=teal!35!black,title={Before you move on}]
What does \ja{もと} mean? (\textquotedblleft{}The origin; the cause\textquotedblright{}.) Say it once more.
\end{tcolorbox}
